\documentclass[10pt,a4paper]{amsart}
\usepackage{amsmath,amssymb,amsthm,graphicx}
% Retained minimal preamble for the Stacks Project chapter more-algebra.tex (More on Algebra).
% The chapter file pinned-original-source.tex carries only \input{preamble} on its first line, so no
% preamble text lives inside the chapter itself.  The Stacks Project uses one single shared preamble for
% every chapter; the snapshot stages it identically as sources/stacks-chow/preamble.tex and
% sources/stacks-etale-cohomology/preamble.tex.  The block below is copied mechanically, line for line,
% from sources/stacks-chow/preamble.tex lines 14-18 (the Xy-pic setup) and 154-217 (the native theorem
% declarations together with the mathematical macros).  Nothing was added, removed or redefined; the
% wrapper class and the amsmath/amssymb/amsthm/graphicx base packages are supplied by the assembler.
% Chapter: more-algebra.tex -- More on Algebra.
% For commutative diagrams we use Xy-pic
\usepackage[all]{xy}

% We use 2cell for 2-commutative diagrams.
\xyoption{2cell}

% Theorem environments.
%
\theoremstyle{plain}
\newtheorem{theorem}[subsection]{Theorem}
\newtheorem{proposition}[subsection]{Proposition}
\newtheorem{lemma}[subsection]{Lemma}

\theoremstyle{definition}
\newtheorem{definition}[subsection]{Definition}
\newtheorem{example}[subsection]{Example}
\newtheorem{exercise}[subsection]{Exercise}
\newtheorem{situation}[subsection]{Situation}

\theoremstyle{remark}
\newtheorem{remark}[subsection]{Remark}
\newtheorem{remarks}[subsection]{Remarks}

\numberwithin{equation}{subsection}

% Macros
%
\def\lim{\mathop{\mathrm{lim}}\nolimits}
\def\colim{\mathop{\mathrm{colim}}\nolimits}
\def\Spec{\mathop{\mathrm{Spec}}}
\def\Hom{\mathop{\mathrm{Hom}}\nolimits}
\def\Ext{\mathop{\mathrm{Ext}}\nolimits}
\def\SheafHom{\mathop{\mathcal{H}\!\mathit{om}}\nolimits}
\def\SheafExt{\mathop{\mathcal{E}\!\mathit{xt}}\nolimits}
\def\Sch{\mathit{Sch}}
\def\Mor{\mathop{\mathrm{Mor}}\nolimits}
\def\Ob{\mathop{\mathrm{Ob}}\nolimits}
\def\Sh{\mathop{\mathit{Sh}}\nolimits}
\def\NL{\mathop{N\!L}\nolimits}
\def\CH{\mathop{\mathrm{CH}}\nolimits}
\def\proetale{{pro\text{-}\acute{e}tale}}
\def\etale{{\acute{e}tale}}
\def\QCoh{\mathit{QCoh}}
\def\Ker{\mathop{\mathrm{Ker}}}
\def\Im{\mathop{\mathrm{Im}}}
\def\Coker{\mathop{\mathrm{Coker}}}
\def\Coim{\mathop{\mathrm{Coim}}}

% Boxtimes
%
\DeclareMathSymbol{\boxtimes}{\mathbin}{AMSa}{"02}

%
% Macros for moduli stacks/spaces
%
\def\QCohstack{\mathcal{QC}\!\mathit{oh}}
\def\Cohstack{\mathcal{C}\!\mathit{oh}}
\def\Spacesstack{\mathcal{S}\!\mathit{paces}}
\def\Quotfunctor{\mathrm{Quot}}
\def\Hilbfunctor{\mathrm{Hilb}}
\def\Curvesstack{\mathcal{C}\!\mathit{urves}}
\def\Polarizedstack{\mathcal{P}\!\mathit{olarized}}
\def\Complexesstack{\mathcal{C}\!\mathit{omplexes}}
% \Pic is the operator that assigns to X its picard group, usage \Pic(X)
% \Picardstack_{X/B} denotes the Picard stack of X over B
% \Picardfunctor_{X/B} denotes the Picard functor of X over B
\def\Pic{\mathop{\mathrm{Pic}}\nolimits}
\def\Picardstack{\mathcal{P}\!\mathit{ic}}
\def\Picardfunctor{\mathrm{Pic}}
\def\Deformationcategory{\mathcal{D}\!\mathit{ef}}

% --- appended 2026-10-12 (Chief): the project own definitions of the reference/slogan/history
% display environments, copied mechanically from the shared Stacks preamble lines 8-12
% (\usepackage{verbatim} plus the three \newenvironment lines, which the chapter preamble keeps
% verbatim from the project preamble). Nothing added, removed or redefined: these environments
% carry no mathematical content and the project itself typesets them as comments. Without these
% three lines a counted statement whose own text uses them fails with Environment slogan undefined.
\usepackage{verbatim}
\newenvironment{reference}{\comment}{\endcomment}
\newenvironment{slogan}{\comment}{\endcomment}
\newenvironment{history}{\comment}{\endcomment}
\usepackage{cleveref}
\crefname{originalsection}{original section}{original sections}
\Crefname{originalsection}{Original section}{Original sections}
\makeatletter
\newlabel{situation-koszul}{{0BKC}{source}{}{sourcepointer-1}{}}
\newlabel{situation-koszul@cref}{{[originalsection][0][]0BKC}{[1][0][]source}{}{}{}}
\makeatother
\crefname{originallemma}{original lemma}{original lemmas}
\Crefname{originallemma}{Original lemma}{Original lemmas}
\makeatletter
\newlabel{lemma-derived-completion}{{091V}{source}{}{sourcepointer-2}{}}
\newlabel{lemma-derived-completion@cref}{{[originallemma][0][]091V}{[1][0][]source}{}{}{}}
\makeatother
\crefname{originallemma}{original lemma}{original lemmas}
\Crefname{originallemma}{Original lemma}{Original lemmas}
\makeatletter
\newlabel{lemma-extended-alternating-Cech-is-colimit-koszul}{{0913}{source}{}{sourcepointer-3}{}}
\newlabel{lemma-extended-alternating-Cech-is-colimit-koszul@cref}{{[originallemma][0][]0913}{[1][0][]source}{}{}{}}
\makeatother
\crefname{originallemma}{original lemma}{original lemmas}
\Crefname{originallemma}{Original lemma}{Original lemmas}
\makeatletter
\newlabel{lemma-dual-perfect-complex}{{07VI}{source}{}{sourcepointer-4}{}}
\newlabel{lemma-dual-perfect-complex@cref}{{[originallemma][0][]07VI}{[1][0][]source}{}{}{}}
\makeatother
\crefname{originallemma}{original lemma}{original lemmas}
\Crefname{originallemma}{Original lemma}{Original lemmas}
\makeatletter
\newlabel{lemma-colim-and-lim-of-duals}{{0BKB}{source}{}{sourcepointer-5}{}}
\newlabel{lemma-colim-and-lim-of-duals@cref}{{[originallemma][0][]0BKB}{[1][0][]source}{}{}{}}
\makeatother
\title{The Koszul derived limit functor is the derived completion adjoint (second proof)}
\author{The Stacks Project Authors}
\date{stacks-project:more-algebra.tex:local-snapshot-20261012; GNU FDL 1.2 or later}
\begin{document}
\maketitle
\setcounter{subsection}{1}
\noindent\textit{Editorial note.} All delivered mathematical text is the author's own, copied from the pinned Stacks Project chapter file more-algebra.tex; no statement or proof step was rewritten, reordered or removed. The counted span is the author environment and its own proof environment. References inside the span are handed to the assembler as navigation stubs that carry the canonical tag of the cited tag only; whether each of them is load-bearing is recorded on the stub itself.\par\medskip
\noindent\textit{Allowed standard prerequisites (editorial).}\begin{itemize}
\item The derived completion adjoint is given by RHom out of the extended alternating Cech complex on f\_1, ..., f\_r. invoked at proof lines 26406-26413. Provides the explicit model of K\textasciicircum{}wedge.
\item The extended alternating Cech complex is the colimit of the Koszul chain complexes K\textasciicircum{}n = K(A, f\_1\textasciicircum{}n, ..., f\_r\textasciicircum{}n), whose duals are the Koszul cochain complexes K\_n. invoked at proof lines 26414-26428. Rewrites the colimit as a hocolim of perfect complexes with controlled duals.
\item For a system of perfect complexes K\_n with colimit K, RHom(hocolim K\textasciicircum{}n, K) = Rlim (K tensor K\_n). invoked at proof lines 26432-26433. Concludes the identification of the two functors.
\end{itemize}\medskip
\ifdefined\hypertarget\hypertarget{sourcepointer-1}{}\fi
\noindent\textit{Source reference.} Original section 0BKC: sources/stacks-more-algebra/pinned-original-source.tex, source lines 26305--26305. Referenced at statement line 26389 for the setting; the situation is used inside the counted span.\par\medskip
\ifdefined\hypertarget\hypertarget{sourcepointer-2}{}\fi
\noindent\textit{Source reference.} Original lemma 091V: sources/stacks-more-algebra/pinned-original-source.tex, source lines 26154--26154. Cited at proof lines 26406-26407 for the explicit model of the adjoint; the conclusion is used inside the counted span.\par\medskip
\ifdefined\hypertarget\hypertarget{sourcepointer-3}{}\fi
\noindent\textit{Source reference.} Original lemma 0913: sources/stacks-more-algebra/pinned-original-source.tex, source lines 7100--7100. Cited at proof lines 26414-26416 to express the Cech complex as a colimit of Koszul chain complexes; the conclusion is used inside the counted span.\par\medskip
\ifdefined\hypertarget\hypertarget{sourcepointer-4}{}\fi
\noindent\textit{Source reference.} Original lemma 07VI: sources/stacks-more-algebra/pinned-original-source.tex, source lines 19176--19176. Cited at proof lines 26424-26426 for the duals of the Koszul complexes; the conclusion is used inside the counted span.\par\medskip
\ifdefined\hypertarget\hypertarget{sourcepointer-5}{}\fi
\noindent\textit{Source reference.} Original lemma 0BKB: sources/stacks-more-algebra/pinned-original-source.tex, source lines 19236--19236. Cited at proof lines 26432-26433 for the final identification of the functors; the conclusion is used inside the counted span.\par\medskip
\begin{lemma}
\label{lemma-derived-completion-koszul}
In Situation \ref{situation-koszul}.
The functor which sends $K \in D(A)$ to the derived limit
$K' = R\lim( K \otimes_A^\mathbf{L} K_n^\bullet )$ is the left
adjoint to the inclusion functor $D_{comp}(A) \to D(A)$
constructed in Lemma \ref{lemma-derived-completion}.
\end{lemma}

\begin{proof}[Second proof]
Denote $K \mapsto K^\wedge$ the adjoint constructed in
Lemma \ref{lemma-derived-completion}. By that lemma we have
$$
K^\wedge = R\Hom\left((A \to \prod\nolimits_{i_0} A_{f_{i_0}} \to
\prod\nolimits_{i_0 < i_1} A_{f_{i_0}f_{i_1}}
\to \ldots \to A_{f_1\ldots f_r}), K\right)
$$
In Lemma \ref{lemma-extended-alternating-Cech-is-colimit-koszul}
we have seen that the extended alternating {\v C}ech complex
$$
A \to \prod\nolimits_{i_0} A_{f_{i_0}} \to
\prod\nolimits_{i_0 < i_1} A_{f_{i_0}f_{i_1}}
\to \ldots \to A_{f_1\ldots f_r}
$$
is a colimit of the Koszul complexes
$K^n = K(A, f_1^n, \ldots, f_r^n)$ sitting in
degrees $0, \ldots, r$. Note that $K^n$ is a finite chain complex
of finite free $A$-modules with dual (as in
Lemma \ref{lemma-dual-perfect-complex})
$R\Hom_A(K^n, A) = K_n$ where $K_n$ is the Koszul cochain
complex sitting in degrees $-r, \ldots, 0$ (as usual).
Thus it suffices to show that
$$
R\Hom_A(\text{hocolim} K^n, K) =  R\lim (K \otimes_A^\mathbf{L} K_n)
$$
This follows from
Lemma \ref{lemma-colim-and-lim-of-duals}.
\end{proof}

\end{document}
